\section{Sharp coefficients and finite controls}\label{sec:sharpness}

The following examples are exact finite entropy computations. They
test the accounting rules, not the existence of a native tape
transposer. Write $\lambda=\log 2$ for the entropy of a fair bit.

\begin{proposition}[Necessary coefficients]\label{prop:coefficients}
\formal{Sharpness.lean}{oneBit_decomposable,residual_bound_attained,residual_coefficient_at_least_one,private_helper_recovers,private_helper_coefficient_at_least_one,shared_helper_coefficient_at_least_two,shared_helper_one_fails}
The following universal coefficients cannot be reduced:
\begin{enumerate}[label=(\roman*)]
\item the coefficient $1$ on the residual sum in
\eqref{eq:conditional-soft};
\item the coefficient $1$ when a single private helper pays for a
single recovered fair bit;
\item the coefficient $2$ on an unrestricted shared helper that
serves the two off-diagonal obligations of the earlier XOR row protocol.
\end{enumerate}
For (iii), at every integer width $b>0$ the protocol is correct on the
full $2\times2$ input cube, has no message cost, and has a helper of
entropy $b\lambda$. Any bound $2b\le c\,H(Y)/\lambda$ for these
obligations therefore requires $c\ge2$.
\end{proposition}
\begin{outline}
For (i), take the empty relation on one coordinate and let that
coordinate be a fair bit. The relation has a certified decomposition;
the two sides of the soft bound are both $2\lambda$. A coefficient
$c$ on the residual would require $2\lambda\le\lambda+c\lambda$.
For (ii), take the private helper to be the bit itself, so exact recovery
costs $\lambda$. For (iii), reuse the full-cube XOR protocol and its
entropy calculation from \cite{SecondPaper2026}: a shared parity of
two independent $b$-bit off-diagonal words lets each row recover the
other word. Its correctness and helper entropy are checked by the
imported \leanfile{TransposeFollowup.SeparatorSharpness} declarations
included in the audit. Substitution gives $2b\le cb$.
\end{outline}
These examples constrain universal capacity coefficients. They do not
preclude savings from the more detailed quantities in
\cref{thm:information}.

\begin{proposition}[Retention is a strict relaxation]\label{prop:retention-control}
\formal{Sharpness.lean}{copied_contexts_retained,copied_contexts_not_nested,xor_uses_not_nested,xor_uses_not_retained}
There is a finite entropy instance and a use list satisfying semantic
retention but not literal nesting. Conversely, the two XOR uses with
contexts $\{B\}$ and $\{A\}$, targeting $A$ and $B$ respectively,
fail semantic retention whenever $H(B\mid A)>0$; they also fail literal
nesting.
\end{proposition}
\begin{outline}
Let $A,B$ be independent fair bits, let $A'$ be a separately named
copy of $A$, and include a helper holding $(A,B)$. The uses
$(\varnothing,A)$ and $(\{A'\},B)$ retain $A$ through its copy, while
the second context does not contain the first target's label.
The formal observations embed these values in a common pair alphabet.
For the XOR list, the second use would have to retain $\{A,B\}$ in
context $\{A\}$; its missing entropy is $H(B\mid A)$.
\end{outline}
Thus a change of representation can preserve a cheap certificate, while
the unrestricted XOR sharing example does not incorrectly receive a
one-capacity retained-list bound.

\begin{proposition}[Loss after a reveal]\label{prop:conditioning-control}
\formal{Sharpness.lean}{parityRank_isPolymatroid,parityRank_singleton,parityRank_of_pair_subset,helperLoss_strictly_smaller,helperLoss_increases_after_reveal}
Let $A,B$ be independent fair bits and let $P=A\oplus B$.
Use helper $A$ twice, each time in empty context with target $B$,
starting with $L_0=\varnothing$. Before any common reveal,
\[
D_A(\varnothing)=0,\qquad \mathcal L(\varnothing)=\lambda.
\]
After revealing $P$,
\[
D_A(\{P\})=\lambda.
\]
Therefore helper-specific loss can be strictly smaller than whole-context
loss, and it need not decrease under additional common conditioning.
\end{proposition}
\begin{outline}
Each of $A,B,P$ has entropy $\lambda$ and every pair has entropy
$2\lambda$; the Lean proof computes these entropies from the uniform
law on two Boolean variables. Before revelation $I(A;B)=0$, so the
transition from $\{B\}$ to $\varnothing$ loses no information about
$A$, although it loses one bit of context. After revelation
$I(A;B\mid P)=\lambda$, so the same transition loses one bit of
helper information.
\end{outline}
Before the reveal, helper $A$ does not decode target $B$; the recovery
residual must be retained. This example establishes a strict comparison
of charges and a failed monotonicity claim, not a cheaper certificate
for a correct transposer. In particular, $D_{\varnothing}$ cannot be
used in place of $D_Q$ after optimizing the choice of $Q$.

\begin{proposition}[A revelation credit]\label{prop:credit-control}
\formal{Sharpness.lean}{oneBitRank_univ,oneBit_reveal_credit}
For a single fair selected bit $X$, take $b_X=\lambda$, no initial
common context, and reveal $X$ itself. Then
\[
2J_Q-S_Q=\lambda,
\qquad 2H(Y_Q)=2\lambda.
\]
Thus the effective revelation charge in
\eqref{eq:information-bound} can be strictly smaller than the charge
in \eqref{eq:entropy-bound}.
\end{proposition}
\begin{outline}
The reveal supplies $J_Q=\lambda$ and leaves zero singleton entropy,
giving $S_Q=\lambda$. The formal one-bit instance evaluates this
difference exactly.
\end{outline}
The credit does not lower the unrestricted shared-helper coefficient:
revealing a parity need not reduce either individual coordinate's
uncertainty. The savings depend on the actual joint law.
